\subsection{Operator Functions}


% **Role:** structure revealed by invariant subspaces and eigenbehavior.
%
% Polynomial functional calculus, matrix/operator functions, exponentials, square roots, spectral mapping.

Operator functions are defined by applying scalar functions to spectral data.\\
Corollary \ref{cor:normal-spectral-resolution-functional-calculus} gives $f(T)=\sum_j f(\mu_j)P_j$ for normal operators; this subsection extends that rule to diagonalizable and Jordan forms, with nilpotent terms where needed.

\begin{theorem} \label{thm:polynomial-functional-calculus-diagonalizable-matrix} \hlt{Polynomial Functional Calculus for Diagonalizable Matrices}\\
Let $A\in\operatorname{Mat}(n,n,\mathbb{F})$ be diagonalizable and write $A=S\Lambda S^{-1}$, where $\Lambda=\operatorname{diag}(\lambda_1,\ldots,\lambda_n)$.\\
If $p\in\mathbb{F}[z]$, then $p(A)=S\operatorname{diag}(p(\lambda_1),\ldots,p(\lambda_n))S^{-1}$.
\begin{proof}
For every $j\geq0$, $A^j=S\Lambda^jS^{-1}$.\\
If $p(z)=\sum_j c_jz^j$, then $p(A)=\sum\limits_j c_jA^j =S\left(\sum\limits_j c_j\Lambda^j\right)S^{-1} =S p(\Lambda) S^{-1}$.\\
Since $\Lambda$ is diagonal, $p(\Lambda)=\operatorname{diag}(p(\lambda_1),\ldots,p(\lambda_n))$.\\
\hlt{Structural Commentary}: Diagonalization reduces $p(A)$ to scalar evaluation on eigenvalues.\\
\hlt{Key Insight}: Compute functions of $A$ by applying the function to the eigenvalues.
\end{proof}
\end{theorem}

\begin{definition} \label{def:normal-spectral-resolution-notation} \hlt{Spectral Resolution Notation}\\
Let $T\in\LL(\mathcal{V})$ be normal on a finite-dimensional complex inner-product space.\\
Define $T=\sum\limits_{\lambda\in\Omega_T}\lambda P_{\lambda}$, where $\Omega_T=\operatorname{spec}(T)$.
\begin{enumerate}[label=(\roman*)]
\setlength{\itemsep}{0pt}
\item \hlt{Spectral Set}: $\Omega_T=\operatorname{spec}(T)$.
\item \hlt{Spectral Projection}: $P_{\lambda}$ is the orthogonal projection onto $E_{\lambda}(T)$.
\item \hlt{Spectral Resolution}: $\{P_{\lambda}\}_{\lambda\in\Omega_T}$, equivalently $T=\sum\limits_{\lambda\in\Omega_T}\lambda P_{\lambda}$.
\end{enumerate}
\end{definition}

\begin{theorem} \label{thm:finite-spectral-functional-calculus-normal} \hlt{Finite Spectral Functional Calculus for Normal Operators}\\
Let $T\in\LL(\mathcal{V})$ be normal, with spectral resolution $T=\sum\limits_{\lambda\in\Omega_T}\lambda P_{\lambda}$.\\
The finite spectral functional calculus is $\mathbb{C}^{\Omega_T}\to\LL(\mathcal{V})$, where $f\mapsto f(T)\defeq\sum\limits_{\lambda\in\Omega_T}f(\lambda)P_{\lambda}$.\\
If $p\in\mathbb{C}[z]$ and $p|_{\Omega_T}=f$, then $p(T)=f(T)$. Hence every $f\in\mathbb{C}^{\Omega_T}$ has a polynomial representative in $T$.
\begin{proof}
For $p|_{\Omega_T}=f$, Corollary \ref{cor:normal-spectral-resolution-functional-calculus} gives $p(T)=\sum\limits_{\lambda\in\Omega_T}p(\lambda)P_{\lambda} =\sum\limits_{\lambda\in\Omega_T}f(\lambda)P_{\lambda} =f(T)$.\\
Since $\Omega_T$ is finite, Lagrange interpolation gives $p\in\mathbb{C}[z]$ with $p|_{\Omega_T}=f$.\\
\hlt{Structural Commentary}: Spectral sums reduce finite functions to polynomial evaluation.\\
\hlt{Key Insight}: Finite spectral data is interpolated by one polynomial.
\end{proof}
\end{theorem}

\begin{theorem} \label{thm:normal-spectral-functional-calculus-rules} \hlt{Rules for Normal Spectral Functions}\\
Let $T\in\LL(\mathcal{V})$ be normal, with spectral resolution $T=\sum\limits_{\lambda\in\Omega_T}\lambda P_{\lambda}$. For $f,g:\Omega_T\to\mathbb{C}$:
\begin{enumerate}[label=(\roman*)]
\setlength{\itemsep}{0pt}
\item $(f+g)(T)=f(T)+g(T)$ and $(fg)(T)=f(T)g(T)$;
\item $f(T)^*=\overline{f}(T)$, where $\overline{f}(\lambda)=\overline{f(\lambda)}$;
\item $f(T)$ is normal and $\operatorname{spec}(f(T))=f(\Omega_T)$.
\end{enumerate}
\begin{proof}
\begin{enumerate}[label=(\roman*)]
\setlength{\itemsep}{0pt}
\item Since $P_{\lambda}P_{\mu}=\delta_{\lambda\mu}P_{\lambda}$, $f(T)g(T)=\sum\limits_{\lambda\in\Omega_T}f(\lambda)g(\lambda)P_{\lambda}=(fg)(T)$; and $(f+g)(T)=f(T)+g(T)$.
\item Since $P_{\lambda}^*=P_{\lambda}$, $f(T)^*=\sum\limits_{\lambda\in\Omega_T}\overline{f(\lambda)}P_{\lambda}=\overline{f}(T)$.
\item For each $\lambda\in\Omega_T$, $f(T)|_{\operatorname{range}(P_{\lambda})}=f(\lambda)I$.\\
Thus $f(T)$ is diagonal in an orthonormal eigenbasis, hence normal, and $\operatorname{spec}(f(T))=f(\Omega_T)$.
\end{enumerate}
\hlt{Structural Commentary}: Spectral projections make operator rules pointwise.\\
\hlt{Key Insight}: Spectral functions act as scalars on spectral subspaces.
\end{proof}
\end{theorem}

\begin{definition} \label{def:positive-operator} \hlt{Positive Operator}\\
$T\in\LL(\mathcal{V})$ is \hlt{positive} if $T=T^*$ and $\innerprod{T\mathbf{v},\mathbf{v}}\geq0$ for all $\mathbf{v}\in\mathcal{V}$.
\end{definition}

\begin{definition} \label{def:square-root-operator} \hlt{Square Root of an Operator}\\
An operator $R\in\LL(\mathcal{V})$ is a \hlt{square root} of $T\in\LL(\mathcal{V})$ if $R^2=T$.
\end{definition}

\begin{theorem} \label{thm:positive-operator-characterization} \hlt{Characterization of Positive Operators}\\
For $T\in\LL(\mathcal{V})$, the following are equivalent:
\begin{enumerate}[label=\roman*.]
\setlength{\itemsep}{0pt}
\item $T$ is positive;
\item $T$ is self-adjoint and every eigenvalue of $T$ is nonnegative;
\item $T$ has a positive square root;
\item $T$ has a self-adjoint square root;
\item there exists $R\in\LL(\mathcal{V})$ such that $T=R^*R$.
\end{enumerate}
\begin{proof}
\begin{enumerate}[label=(\roman*).]
\setlength{\itemsep}{0pt}
\item[$(i)$]$\Rightarrow (ii)$ If $T\mathbf{v}=\lambda\mathbf{v}$ and $\mathbf{v}\neq\mathbf{0}$, then $\lambda\in\mathbb{R}$ by Theorem \ref{thm:self-adjoint-eigenvalues-real}, and $0\leq\innerprod{T\mathbf{v},\mathbf{v}}=\lambda\norm{\mathbf{v}}^2$, so $\lambda\geq0$.
\item[$(ii)$]$\Rightarrow (iii)$ By the real or complex spectral theorem, choose an orthonormal eigenbasis $(\mathbf{e}_1,\ldots,\mathbf{e}_n)$ with $T\mathbf{e}_j=\lambda_j\mathbf{e}_j$ and $\lambda_j\geq0$. Define $S\mathbf{e}_j=\sqrt{\lambda_j}\mathbf{e}_j$ and extend linearly.\\
Then $S$ is diagonal with nonnegative real diagonal entries in an orthonormal basis, hence self-adjoint and positive, and $S^2\mathbf{e}_j=\lambda_j\mathbf{e}_j=T\mathbf{e}_j$ for every $j$. Thus $S^2=T$.
\item[$(iii)$]$\Rightarrow (iv)$ A positive square root is self-adjoint.
\item[$(iv)$]$\Rightarrow (v)$ If $T=R^2$ with $R=R^*$, then $T=R^*R$.
\item[$(v)$]$\Rightarrow (i)$ If $T=R^*R$, then $T^*=T$ and $\innerprod{T\mathbf{v},\mathbf{v}}=\innerprod{R\mathbf{v},R\mathbf{v}}=\norm{R\mathbf{v}}^2\geq0$ for every $\mathbf{v}$, so $T$ is positive.
\end{enumerate}
\hlt{Structural Commentary}: Positivity equals nonnegative spectrum and square-root factorization.\\
\hlt{Key Insight}: In an eigenbasis, the proof is scalar: $\lambda\geq0$ iff $\sqrt{\lambda}$ exists.
\end{proof}
\end{theorem}

\begin{proposition} \label{prop:positive-square-root-unique} \hlt{Unique Positive Square Root}\\
Every positive operator on $\mathcal{V}$ has a unique positive square root.
\begin{proof}
Existence follows from Theorem \ref{thm:positive-operator-characterization}. For uniqueness, decompose $\mathcal{V}=E_{\lambda_1}(T)\oplus\cdots\oplus E_{\lambda_m}(T)$.\\
Let $R$ be a positive square root of $T$. Since $R^2=T$, we have $RT=R^3=TR$, so each $E_{\lambda_j}(T)$ is invariant under $R$. On $E_{\lambda_j}(T)$, the restriction $R_j\defeq R|_{E_{\lambda_j}(T)}$ is positive and satisfies $R_j^2=\lambda_jI$.\\
By the spectral theorem for the positive operator $R_j$, every eigenvalue $\rho$ of $R_j$ satisfies $\rho\geq0$ and $\rho^2=\lambda_j$, hence $\rho=\sqrt{\lambda_j}$. Therefore $R_j=\sqrt{\lambda_j}I$ on $E_{\lambda_j}(T)$.\\
Thus every positive square root agrees with the same scalar operator on each eigenspace $E_{\lambda_j}(T)$, so it is unique.\\
\hlt{Structural Commentary}: Each eigenspace determines its own scalar root.\\
\hlt{Key Insight}: Positivity selects the nonnegative root.
\end{proof}
\end{proposition}

\begin{definition} \label{def:positive-square-root-notation} \hlt{Positive Square Root Notation}\\
If $T$ is positive, then $\sqrt{T}$ denotes the unique positive square root of $T$.
\end{definition}

\begin{example} \label{ex:spectral-functional-calculus-basic-functions} \hlt{Basic Functions from the Spectral Calculus}\\
Let $T\in\LL(\mathcal{V})$ be normal, with spectral resolution $T=\sum\limits_{\lambda\in\Omega_T}\lambda P_{\lambda}$. Then $T^*=\sum\limits_{\lambda\in\Omega_T}\overline{\lambda}P_{\lambda}$.\\
If $0\notin\Omega_T$, then $T^{-1}=\sum\limits_{\lambda\in\Omega_T}\lambda^{-1}P_{\lambda}$. If $T\geq0$, then $T^{1/2}=\sum\limits_{\lambda\in\Omega_T}\lambda^{1/2}P_{\lambda}$, where $\lambda^{1/2}\geq0$.
\end{example}

\begin{theorem} \label{thm:normal-spectral-projection-commutativity} \hlt{Commutativity via Spectral Projections}\\
Let $T,S\in\LL(\mathcal{V})$ be normal, with spectral resolutions $T=\sum\limits_{\lambda\in\Omega_T}\lambda P_{\lambda}$ and $S=\sum\limits_{\mu\in\Omega_S}\mu Q_{\mu}$.
\begin{enumerate}[label=(\roman*)]
\setlength{\itemsep}{0pt}
\item For $R\in\LL(\mathcal{V})$, $RT=TR$ if and only if $RP_{\lambda}=P_{\lambda}R$ for all $\lambda\in\Omega_T$.
\item $TS=ST$ if and only if $P_{\lambda}Q_{\mu}=Q_{\mu}P_{\lambda}$ for all $\lambda\in\Omega_T$ and $\mu\in\Omega_S$.
\item If $f:\Omega_T\to\mathbb{C}$ and $g:\Omega_S\to\mathbb{C}$ are injective, then $f(T)g(S)=g(S)f(T)$ if and only if $TS=ST$.
\end{enumerate}
\begin{proof}
\begin{enumerate}[label=(\roman*).]
\setlength{\itemsep}{0pt}
\item If $RT=TR$, then $Rp(T)=p(T)R$ for all $p\in\mathbb{C}[z]$. Since $P_{\lambda}\in\mathbb{C}[T]$ by Corollary \ref{cor:normal-spectral-resolution-functional-calculus}, $RP_{\lambda}=P_{\lambda}R$. Conversely, if $RP_{\lambda}=P_{\lambda}R$ for all $\lambda\in\Omega_T$, then $RT=\sum\limits_{\lambda\in\Omega_T}\lambda RP_{\lambda} =\sum\limits_{\lambda\in\Omega_T}\lambda P_{\lambda}R =TR$.
\item If $TS=ST$, then (i) with $R=S$ gives $SP_{\lambda}=P_{\lambda}S$. Since $Q_{\mu}\in\mathbb{C}[S]$, $P_{\lambda}Q_{\mu}=Q_{\mu}P_{\lambda}$ for all $\lambda,\mu$. Conversely, if $P_{\lambda}Q_{\mu}=Q_{\mu}P_{\lambda}$ for all $\lambda,\mu$, then $TS=\sum\limits_{\lambda,\mu}\lambda\mu P_{\lambda}Q_{\mu} =\sum\limits_{\lambda,\mu}\lambda\mu Q_{\mu}P_{\lambda} =ST$
\item Suppose $f$ and $g$ are injective. If $f(T)g(S)=g(S)f(T)$, choose $h:f(\Omega_T)\to\Omega_T$ and $k:g(\Omega_S)\to\Omega_S$ with $h\circ f=\operatorname{id}_{\Omega_T}$ and $k\circ g=\operatorname{id}_{\Omega_S}$. By Theorem \ref{thm:finite-spectral-functional-calculus-normal}, $T=h(f(T))$ and $S=k(g(S))$, so $TS=ST$. Conversely, if $TS=ST$, then (ii) gives $P_{\lambda}Q_{\mu}=Q_{\mu}P_{\lambda}$,\\
and $f(T)g(S)=\sum\limits_{\lambda,\mu}f(\lambda)g(\mu)P_{\lambda}Q_{\mu} =\sum\limits_{\lambda,\mu}f(\lambda)g(\mu)Q_{\mu}P_{\lambda} =g(S)f(T)$.
\end{enumerate}
\hlt{Structural Commentary}: Spectral projections are the coordinate idempotents of a normal operator.\\
\hlt{Key Insight}: To commute with a normal operator is to preserve each of its eigenspaces.
\end{proof}
\end{theorem}

\begin{corollary} \label{cor:commuting-normal-pair-generated-by-one-normal-operator} \hlt{A Commuting Normal Pair Comes from One Normal Operator}\\
Let $T,S\in\LL(\mathcal{V})$ be normal on a finite-dimensional complex inner-product space.\\
Then $TS=ST$ $\Longleftrightarrow$ $\exists\, H\in\LL(\mathcal{V})$ normal, $\ p,q\in\mathbb{C}[z]$ such that $T=p(H),\ S=q(H)$.\\
If $TS=ST$, then $H$ may be chosen in $\mathbb{C}[T,S]$.
\begin{proof}
\begin{enumerate}[label=\roman*.]
\setlength{\itemsep}{0pt}
\item[$(\Rightarrow)$] Suppose $TS=ST$. Write $T=\sum\limits_{\lambda\in\Omega_T}\lambda P_{\lambda}$ and $S=\sum\limits_{\mu\in\Omega_S}\mu Q_{\mu}$.\\
By Theorem \ref{thm:normal-spectral-projection-commutativity}, $P_{\lambda}Q_{\mu}=Q_{\mu}P_{\lambda}$. Let $\mathcal{I}=\{(\lambda,\mu)\in\Omega_T\times\Omega_S:P_{\lambda}Q_{\mu}\neq0\}$.\\
The projections $P_{\lambda}Q_{\mu}$, $(\lambda,\mu)\in\mathcal{I}$, are pairwise orthogonal and sum to $I_{\mathcal{V}}$.\\
Choose distinct scalars $\alpha_{\lambda\mu}$ for $(\lambda,\mu)\in\mathcal{I}$ and set $H=\sum\limits_{(\lambda,\mu)\in\mathcal{I}}\alpha_{\lambda\mu}P_{\lambda}Q_{\mu}$.\\
Then $H$ is normal, and $H\in\mathbb{C}[T,S]$ since $P_{\lambda}\in\mathbb{C}[T]$ and $Q_{\mu}\in\mathbb{C}[S]$.\\
Choose $p,q\in\mathbb{C}[z]$ with $p(\alpha_{\lambda\mu})=\lambda$ and $q(\alpha_{\lambda\mu})=\mu$ on $\mathcal{I}$.\\
Hence $p(H)=\sum\limits_{(\lambda,\mu)\in\mathcal{I}}\lambda P_{\lambda}Q_{\mu}=T$, and $q(H)=\sum\limits_{(\lambda,\mu)\in\mathcal{I}}\mu P_{\lambda}Q_{\mu}=S$.
\item[$(\Leftarrow)$] If $T=p(H)$ and $S=q(H)$, then $TS=p(H)q(H)=q(H)p(H)=ST$.
\end{enumerate}
\hlt{Structural Commentary}: One scalar labels each common spectral block.\\
\hlt{Key Insight}: Joint spectral projections reduce a pair of commuting normal operators to finite scalar data.
\end{proof}
\end{corollary}

\begin{definition} \label{def:matrix-exponential} \hlt{Matrix Exponential}\\
For $A\in\operatorname{Mat}(n,n,\mathbb{C})$ and $t\in\mathbb{C}$, define $e^{tA}\defeq\sum\limits_{k=0}^{\infty}\frac{t^kA^k}{k!}$.
\end{definition}

\begin{proposition} \label{prop:diagonalizable-pure-modes-powers-exponentials} \hlt{Pure Modes for Powers and Exponentials}\\
Let $A\in\operatorname{Mat}(n,n,\mathbb{C})$ be diagonalizable with $A=S\Lambda S^{-1}$, where $\Lambda=\operatorname{diag}(\lambda_1,\ldots,\lambda_n)$ and the columns of $S$ are $\mathbf{v}_1,\ldots,\mathbf{v}_n$. If $\mathbf{x}_0=\sum\limits_{j=1}^n c_j\mathbf{v}_j$, then $A^k\mathbf{x}_0=\sum\limits_{j=1}^n c_j\lambda_j^k\mathbf{v}_j$ for $k\geq0$; and $e^{tA}\mathbf{x}_0=\sum\limits_{j=1}^n c_j e^{\lambda_jt}\mathbf{v}_j$ for $t\in\mathbb{C}$.
\begin{proof}
Let $c=(c_1,\ldots,c_n)^t$, so $\mathbf{x}_0=Sc$.\\
Then $A^k=S\Lambda^kS^{-1}$, and Definition \ref{def:matrix-exponential} gives $e^{tA}=S\left(\sum\limits_{r=0}^{\infty}\frac{t^r\Lambda^r}{r!}\right)S^{-1} =S\operatorname{diag}(e^{\lambda_1t},\ldots,e^{\lambda_nt})S^{-1}$.\\
Applying these two formulas to $Sc$ gives the stated identities.\\
\hlt{Structural Commentary}: Eigenbasis coordinates decouple powers and exponentials.\\
\hlt{Key Insight}: Diagonalization turns coupled evolution into independent scalar modes.
\end{proof}
\end{proposition}

\begin{corollary} \label{cor:commutant-distinct-eigenvalues-polynomials} \hlt{Simple Spectrum Commutants Are Polynomials}\\
If $A\in\operatorname{Mat}(n,n,\mathbb{C})$ has $n$ distinct eigenvalues, then $\operatorname{Comm}(A)=\{p(A):p\in\mathbb{C}[z]\}$, $\dim\operatorname{Comm}(A)=n$.
\begin{proof}
Write $A=S\operatorname{diag}(\lambda_1,\ldots,\lambda_n)S^{-1}$ with distinct $\lambda_i$. Theorem \ref{thm:commutant-of-diagonalizable-matrix} says every $B\in\operatorname{Comm}(A)$ has the form $S\operatorname{diag}(\xi_1,\ldots,\xi_n)S^{-1}$. Choose the Lagrange interpolation polynomial $p(z)=\sum\limits_{i=1}^{n}\xi_i\prod\limits_{j\neq i}\frac{z-\lambda_j}{\lambda_i-\lambda_j}$.\\
Then $p(\lambda_i)=\xi_i$ for every $i$, so Theorem \ref{thm:polynomial-functional-calculus-diagonalizable-matrix} gives $B=p(A)$.\\
The dimension statement is the case $n_1=\cdots=n_n=1$ of Theorem \ref{thm:commutant-of-diagonalizable-matrix}.\\
\hlt{Structural Commentary}: Simple spectrum leaves one scalar choice per eigendirection.\\
\hlt{Key Insight}: Interpolation converts arbitrary diagonal data into a polynomial in $A$.
\end{proof}
\end{corollary}

\begin{corollary} \label{cor:polynomial-reduction-by-characteristic-polynomial} \hlt{Reduction Modulo the Characteristic Polynomial}\\
Let $A\in\operatorname{Mat}(n,n,\mathbb{C})$ and let $f\in\mathbb{C}[z]$. If $f=p_Aq+r$ with $\deg r<n$, then $f(A)=r(A)$.
\begin{proof}
Evaluate division identity at $A$: $f(A)=p_A(A)q(A)+r(A)=0q(A)+r(A)=r(A)$.\\
Theorem \ref{thm:cayley-hamilton-matrix} gives $p_A(A)=0$.\\
\hlt{Structural Commentary}: Cayley--Hamilton makes polynomial expressions in $A$ reducible modulo $p_A$.\\
\hlt{Key Insight}: Every high power of $A$ can be rewritten using powers below $n$.
\end{proof}
\end{corollary}

\begin{corollary} \label{cor:powers-of-matrix-span-dimension-at-most-n} \hlt{Powers Span at Most $n$ Dimensions}\\
Let $A\in\operatorname{Mat}(n,n,\mathbb{C})$. Then $\operatorname{span}(I_n,A,A^2,\ldots)=\operatorname{span}(I_n,A,\ldots,A^{n-1})$, so this has dimension at most $n$.
\begin{proof}
For each $k\geq0$, Corollary \ref{cor:polynomial-reduction-by-characteristic-polynomial} applied to $f(z)=z^{n+k}$ gives a polynomial $r_k$ of degree at most $n-1$ such that $A^{n+k}=r_k(A)$. Thus every power of $A$ lies in $\operatorname{span}(I_n,A,\ldots,A^{n-1})$.\\
\hlt{Structural Commentary}: Polynomial reduction collapses the infinite power list to its first $n$ terms.\\
\hlt{Key Insight}: Cayley--Hamilton gives the recurrence relation for powers of $A$.
\end{proof}
\end{corollary}

\begin{corollary} \label{cor:inverse-from-characteristic-polynomial} \hlt{Inverse from the Characteristic Polynomial}\\
Let $A\in\operatorname{Mat}(n,n,\mathbb{C})$ have $p_A(z)=z^n+c_{n-1}z^{n-1}+\cdots+c_1z+c_0$.\\
If $A$ is invertible, then $A^{-1}=-c_0^{-1}(A^{n-1}+c_{n-1}A^{n-2}+\cdots+c_2A+c_1I_n)$.
\begin{proof}
Since $A$ is invertible, $0\notin\operatorname{spec}(A)$, so $c_0=p_A(0)=(-1)^n\det A\neq0$.\\
Cayley--Hamilton gives $A^n+c_{n-1}A^{n-1}+\cdots+c_1A+c_0I_n=0$.\\
Move $c_0I_n$ to the other side and multiply by $-c_0^{-1}A^{-1}$.\\
\hlt{Structural Commentary}: The characteristic equation contains an identity term exactly when $A$ is invertible.\\
\hlt{Key Insight}: Solve Cayley--Hamilton for $A^{-1}$ after factoring out one $A$.
\end{proof}
\end{corollary}

\begin{proposition} \label{prop:eigenvalues-of-polynomial-in-matrix} \hlt{Eigenvalues of a Polynomial in a Matrix}\\
Let $A\in\operatorname{Mat}(n,n,\mathbb{C})$, let $p\in\mathbb{C}[z]$, and let $\lambda_1,\ldots,\lambda_n$ be the eigenvalues of $A$, counted with algebraic multiplicity. Then the eigenvalues of $p(A)$ are $p(\lambda_1),\ldots,p(\lambda_n)$, counted with algebraic multiplicity.
\begin{proof}
Choose a Schur form $U^*AU=T$ with diagonal $\lambda_1,\ldots,\lambda_n$. Then $p(A)=Up(T)U^*$, and $p(T)$ is upper triangular with diagonal $p(\lambda_1),\ldots,p(\lambda_n)$. Hence $p(A)$ has precisely these eigenvalues.\\
\hlt{Structural Commentary}: Schur form reduces the claim to the diagonal.\\
\hlt{Key Insight}: A polynomial preserves triangular form and applies $p$ to each diagonal entry.
\end{proof}
\end{proposition}

\begin{corollary} \label{cor:polynomial-spectral-mapping-matrix} \hlt{Polynomial Spectral Mapping}\\
Let $A\in\operatorname{Mat}(n,n,\mathbb{C})$ and $p\in\mathbb{C}[z]$. Then $\operatorname{spec}(p(A))=\{p(\lambda):\lambda\in\operatorname{spec}(A)\}$.
\begin{proof}
Let $\lambda_1,\ldots,\lambda_n$ be the eigenvalues of $A$, counted with algebraic multiplicity.\\
By Proposition \ref{prop:eigenvalues-of-polynomial-in-matrix}, the eigenvalues of $p(A)$, counted with algebraic multiplicity, are $p(\lambda_1),\ldots,p(\lambda_n)$. Taking underlying sets gives the claim.\\
\hlt{Structural Commentary}: Polynomial evaluation pushes the spectrum forward through the scalar polynomial.\\
\hlt{Key Insight}: Once Schur form has exposed the diagonal, $p$ acts entry by entry there.
\end{proof}
\end{corollary}

\begin{corollary} \label{cor:polynomial-spectral-mapping-operator} \hlt{Polynomial Spectral Mapping for Operators}\\
Let $\mathcal{V}\neq\{\mathbf{0}\}$ be finite-dimensional over $\mathbb{C}$, let $T\in\LL(\mathcal{V})$, and let $p\in\mathbb{C}[z]$.\\
Then $\operatorname{spec}(p(T))=\{p(\lambda):\lambda\in\operatorname{spec}(T)\}$.
\begin{proof}
Choose a basis $\beta$ and set $A=[T]_{\beta\leftarrow\beta}$.\\
Then $[p(T)]_{\beta\leftarrow\beta}=p(A)$, and coordinate isomorphism gives $\operatorname{spec}(T)=\operatorname{spec}(A)$ and $\operatorname{spec}(p(T))=\operatorname{spec}(p(A))$.\\
Corollary \ref{cor:polynomial-spectral-mapping-matrix} gives $\operatorname{spec}(p(T))=\operatorname{spec}(p(A))=\{p(\lambda):\lambda\in\operatorname{spec}(A)\}=\{p(\lambda):\lambda\in\operatorname{spec}(T)\}$.\\
\hlt{Structural Commentary}: Coordinates reduce the operator statement to the matrix statement.\\
\hlt{Key Insight}: Polynomial spectral mapping is an operator fact seen through any convenient basis.
\end{proof}
\end{corollary}

\begin{corollary} \label{cor:polynomial-in-matrix-invertible-spectrum-test} \hlt{Polynomial Invertibility Test}\\
Let $A\in\operatorname{Mat}(n,n,\mathbb{C})$ and $p\in\mathbb{C}[z]$. If $p(\lambda)\neq0$ for every $\lambda\in\operatorname{spec}(A)$, then $p(A)$ is invertible.
\begin{proof}
By Proposition \ref{prop:eigenvalues-of-polynomial-in-matrix}, $0\notin\operatorname{spec}(p(A))$. Corollary \ref{cor:matrix-invertible-iff-zero-not-spectrum} gives invertibility.\\
\hlt{Structural Commentary}: Invertibility of $p(A)$ is checked on spectral values of $A$.\\
\hlt{Key Insight}: Avoiding zero on the spectrum prevents $p(A)$ from having a kernel.
\end{proof}
\end{corollary}

\begin{proposition} \label{prop:exponential-of-jordan-block} \hlt{Exponential of a Jordan Block}\\
For $k\geq1$, $\lambda\in\mathbb{C}$, and $t\in\mathbb{C}$, $e^{tJ_k(\lambda)} =e^{\lambda t}\sum\limits_{r=0}^{k-1}\frac{t^r}{r!}J_k(0)^r$.
\begin{proof}
Since $J_k(\lambda)=\lambda I_k+J_k(0)$ and the two summands commute, $e^{tJ_k(\lambda)}=e^{\lambda tI_k}e^{tJ_k(0)}$.\\
The first factor is $e^{\lambda t}I_k$, and the second exponential truncates because $J_k(0)^k=0$.\\
\hlt{Structural Commentary}: The scalar part supplies $e^{\lambda t}$; the nilpotent part supplies a finite polynomial in $t$.\\
\hlt{Key Insight}: Nilpotence turns the exponential series into a finite sum.
\end{proof}
\end{proposition}

\begin{lemma} \label{lem:identity-plus-nilpotent-square-root} \hlt{Identity Plus Nilpotent Has a Square Root}\\
Let $\mathcal{V}$ be a complex vector space and let $N\in\LL(\mathcal{V})$ be nilpotent. Then $I_{\mathcal{V}}+N$ has a square root.
\begin{proof}
Choose $m\geq1$ such that $N^m=0$. If $m=1$, then $N=0$ and $I_{\mathcal{V}}$ is a square root.\\
Otherwise seek $R=I_{\mathcal{V}}+a_1N+\cdots+a_{m-1}N^{m-1}$. With $a_0=1$, choose $a_1,\ldots,a_{m-1}$ recursively so the coefficients of $R^2$ match those of $I_{\mathcal{V}}+N$: $2a_1=1$, and $2a_k+\sum\limits_{r=1}^{k-1}a_ra_{k-r}=0\quad (2\leq k\leq m-1)$\\
Terms of degree $\geq m$ vanish because $N^m=0$, so $R^2=I_{\mathcal{V}}+N$.\\
\hlt{Structural Commentary}: Nilpotence reduces the square-root series to a finite polynomial.\\
\hlt{Key Insight}: Match coefficients recursively.
\end{proof}
\end{lemma}

\begin{figure}[H]
  \centering
  \begin{tikzpicture}[
  >=Stealth,
  every node/.style={font=\small},
  box/.style={draw=gray!60, rounded corners=3pt, fill=gray!5, inner sep=5pt, align=center},
  arrow/.style={->, line width=0.85pt, blue!70!black}
]
  \node[box] (nil) at (0,0) {$N^m=0$};
  \node[box] (poly) at (3.5,0) {$R=I+a_1N+\cdots+a_{m-1}N^{m-1}$};
  \node[box] (sq) at (7.0,0) {$R^2=I+N$};

  \draw[arrow] (nil) -- (poly);
  \draw[arrow] (poly) -- (sq);
\end{tikzpicture}

  \caption{Nilpotence truncates the square-root series; coefficients are chosen so $R^2=I+N$.}
  \label{fig:nilpotent-square-root-polynomial}
\end{figure}

\begin{theorem} \label{thm:complex-invertible-operator-has-square-root} \hlt{Invertible Complex Operators Have Square Roots}\\
Let $\mathcal{V}$ be finite-dimensional over $\mathbb{C}$ and let $T\in\LL(\mathcal{V})$ be invertible. Then $T$ has a square root.
\begin{proof}
Let $\mathcal{U}_j=\mathcal{G}_{\lambda_j}(T)$ for the distinct eigenvalues of $T$. Primary decomposition gives $\mathcal{V}=\mathcal{U}_1\oplus\cdots\oplus\mathcal{U}_m$, and invertibility gives $\lambda_j\neq0$. Put $N_j=(T-\lambda_jI_{\mathcal{V}})|_{\mathcal{U}_j}$.\\
Then $T|_{\mathcal{U}_j} =\lambda_j\left(I_{\mathcal{U}_j}+\lambda_j^{-1}N_j\right)$ with $\lambda_j^{-1}N_j$ nilpotent. By Lemma \ref{lem:identity-plus-nilpotent-square-root}, choose $B_j^2=I_{\mathcal{U}_j}+\lambda_j^{-1}N_j$; choose $\alpha_j^2=\lambda_j$ and set $R_j=\alpha_jB_j$. Then $R_j^2=T|_{\mathcal{U}_j}$. Define $R=R_1\oplus\cdots\oplus R_m$ on the direct sum. Then $R^2=T$.\\
\hlt{Structural Commentary}: Take square roots on primary blocks, then recombine.\\
\hlt{Key Insight}: Invertibility makes every scalar block nonzero.
\end{proof}
\end{theorem}
